\documentclass[11pt,a4paper]{article}
\usepackage{curso}
\begin{document}
\portada{Lección}{5.2}{p-bits: MTJ y CMOS}

\noindent\textbf{Unidad 5.}\quad\textbf{Objetivos:} pbit físico con una unión túnel magnética de barrera baja; pbit digital (RNG + LUT + comparador);
p-computadores y máquinas de Ising dispersas.
\textbf{Prerrequisitos:} 4.5 (pbit), 3.7 (retención de Arrhenius).

\section{La unión túnel magnética estocástica}
Dos capas ferromagnéticas separadas por un aislante: resistencia baja si las magnetizaciones son paralelas, alta si son antiparalelas. La capa libre
tiene una barrera $\Delta$ entre ambas orientaciones y salta con el tiempo medio de Néel--Arrhenius
\begin{equation}
\tau=\tau_0\,e^{\Delta/\kB T},\qquad\tau_0\sim1\ \text{ns}.
\end{equation}
\begin{center}\begin{tabular}{lccccc}\toprule
$\Delta/\kB T$ & 1 & 5 & 10 & 20 & 40\\\midrule
$\tau$ & 2,7 ns & 0,15 $\mu$s & 22 $\mu$s & 0,49 s & $\approx7{,}5$ años\\\bottomrule\end{tabular}\end{center}
Con $\Delta\approx40\,\kB T$ es una memoria; con $\Delta\sim\kB T$ el ruido térmico la voltea en nanosegundos: un pbit. Una corriente o voltaje que inclina el
pozo doble ajusta la fracción de tiempo en cada estado según una sigmoide; en el circuito típico 1T/1MTJ + inversor la entrada es el voltaje de compuerta
\cite{camsari}.
\begin{intuicion}
El mismo dispositivo es memoria o generador aleatorio según la altura de la barrera. El pbit aprovecha lo que la memoria combate.
\end{intuicion}

\section{pbits digitales}
$s=\mathrm{sgn}(\mathrm{LUT}_{\tanh}(u)-r)$ con $r$ de un PRNG (p.~ej. LFSR): la regla de la lección 4.5 en lógica estándar. No ahorra la energía de la
aleatoriedad, pero se construye hoy en FPGA con $\sim10^5$ pbits y conectividad arbitraria: es la plataforma de desarrollo de los p-computadores.

\section{p-computadores y máquinas de Ising dispersas}
Muchos pbits conectados según un grafo y actualizados con Gibbs cromático (lección 4.3). Con grado bajo, pocos colores: en cada fase de reloj se actualizan
todos los pbits de un color. Las máquinas de Ising dispersas en FPGA alcanzan miles de millones de volteos por segundo \cite{aadit}.

\begin{aplicacion}
Optimización, muestreo y lógica invertible (lección 4.10). Retos: variabilidad entre dispositivos y calibración del sesgo de cada pbit. Laboratorio L8:
circuitos de pbits simulados, curva de transferencia y tiempo de correlación.
\end{aplicacion}

\begin{error}
\begin{itemize}[nosep]
\item Usar una MTJ de memoria ($\Delta\approx40\,\kB T$) como pbit. \quad Ignorar la variabilidad: cada pbit necesita calibración.
\item Actualizar vecinos en la misma fase. \quad Contar volteos/s sin verificar la distribución muestreada.
\end{itemize}
\end{error}

\begin{thebibliography}{9}\small
\bibitem[17]{camsari} K. Y. Camsari, B. M. Sutton y S. Datta, p-bits for probabilistic spin logic, Appl. Phys. Rev. 6, 011305 (2019).
\bibitem[18]{kaiser} J. Kaiser y S. Datta, Probabilistic computing with p-bits, Appl. Phys. Lett. 119, 150503 (2021).
\bibitem[19]{aadit} N. A. Aadit et al., Massively parallel probabilistic computing with sparse Ising machines, Nat. Electron. 5, 460 (2022).
\end{thebibliography}
\end{document}
